\section{Normal estimation on digital surfaces}

In this section, we write $\widehat{i}$ for the $d-1$ subset of indices
$\{1,\ldots,d\} \setminus \{ i \}$.

\begin{theorem}
\label{thm:normal-recovery}
Let $T\subset\mathbb R^d$ be a $(d-1)$-dimensional linear subspace
with unit normal vector $n=(n_1,\ldots,n_d)$.
Then \( V_{\widehat{i}}(T,p,r)=\omega_{d-1}r^{d-1}|n_i|. \)
Consequently, the $d$ $\pi_{\widehat{i}}$-projection volumes uniquely
determine the absolute values of the components
of the unit normal vector:
\[
W_{d-1}(T,p,r) := \sum_{i=1}^d V_{\widehat{i}}(T,p,r)^2
=
\omega_{d-1}^2r^{2(d-1)},
~\text{and}~
|n_i|=\frac{V_{\widehat{i}}(T,p,r)}
{\sqrt{W_{d-1}(T,p,r)}}.
\]
\end{theorem}

\begin{proof}
Let $A\in\mathbb R^{d\times(d-1)}$ have orthonormal columns spanning
$T$, so that $A^\top A=\mathrm{Id}_{d-1}$.  For each $i, 1\Le i \Le
d$, let $A_{\widehat i}$ be the square matrix obtained by deleting the
$i$-th row of $A$.  The restriction of $\pi_{\widehat i}$ to $T$ is
represented by $A_{\widehat i}$.  Hence, by standard area computation, 
\[
V_{\widehat{i}}(T,p,r)=\omega_{d-1}r^{d-1}|\det A_{\widehat i}|.
\]

\begin{sloppypar}
Consider the vector of signed maximal minors $c=(c_1,\ldots,c_d)$,
with $c_i = (-1)^{i+1} \det A_{\widehat i}$.  By the cofactor expansion,
$c$ is orthogonal to every column of $A$, and therefore $c$ is
parallel to $n$.
Moreover, applying the Cauchy--Binet formula gives
\(
\|c\|_2^2
=
\sum_{i=1}^d\det(A_{\widehat i})^2
=
\det(A^\top A)=1.
\).
Since $n$ is a unit normal vector and $c$ is parallel to $n$, we obtain
$c=\pm n$. In particular,
\(
|\det A_{\widehat i}|=|n_i|.
\)
Substitution into the expression for $V_{\widehat{i}}(T,p,r)$
proves the first identity. The remaining
identities follow from $\|n\|_2=1$. \qed
\end{sloppypar}
\end{proof}



\begin{corollary}[Local estimation of the absolute normal components]
\label{cor:local-normal-estimation}
Let $X\subset\mathbb R^d$ be a $C^2$ hypersurface
without boundary, with reach at least $\rho>0$.
Let $p\in X$, and let
$n(p)=(n_1(p),\ldots,n_d(p))$
be a unit normal vector at $p$.

For each $i\in\{1,\ldots,d\}$, let
$\pi_{\widehat i}:\mathbb R^d\to\mathbb R^{d-1}$
denote the orthogonal coordinate projection
obtained by deleting the $i$-th coordinate.
Define
\[
V_i(p,r)
:=
\mathcal H^{d-1}
\bigl(\pi_{\widehat i}(X\cap B(p,r))\bigr),
\]
and the normalized estimator
\[
\widehat n_i(p,r)
:=
\frac{V_i(p,r)}
{\left(\sum_{j=1}^d V_j(p,r)^2\right)^{1/2}}.
\]

There exist constants $c_d,C_d>0$, depending
only on $d$, such that, whenever
$0<r<c_d\rho$, the estimator is well-defined
and satisfies
\[
\boxed{
\left\|
\widehat n(p,r)-|n(p)|
\right\|_2
\le C_d\frac r\rho,
}
\]
where
\[
|n(p)|=(|n_1(p)|,\ldots,|n_d(p)|).
\]
\end{corollary}

\begin{proof}
Let $T=T_pX$ and let $B_T(p,r)$ denote the
$(d-1)$-dimensional ball of radius $r$
contained in the affine tangent hyperplane
$p+T$.

By the local stability theorem for coordinate
projection volumes, there exists a constant
$C_d^{\mathrm{stab}}>0$ such that, for every $i$,
\[
\left|
V_i(p,r)
-
\mathcal H^{d-1}
\bigl(\pi_{\widehat i}(B_T(p,r))\bigr)
\right|
\le
C_d^{\mathrm{stab}}\frac{r^d}{\rho}.
\]

By the exact formula for the coordinate
projections of a tangent ball,
\[
\mathcal H^{d-1}
\bigl(\pi_{\widehat i}(B_T(p,r))\bigr)
=
\omega_{d-1}r^{d-1}|n_i(p)|.
\]
Consequently,
\[
\left|
V_i(p,r)
-
\omega_{d-1}r^{d-1}|n_i(p)|
\right|
\le
C_d^{\mathrm{stab}}\frac{r^d}{\rho}.
\]

Set
\[
a=\omega_{d-1}r^{d-1},
\qquad
v=(V_1(p,r),\ldots,V_d(p,r)),
\qquad
u=|n(p)|.
\]
Since $\|u\|_2=1$, the preceding estimates give
\[
\|v-au\|_2
\le
\sqrt d\,C_d^{\mathrm{stab}}\frac{r^d}{\rho}.
\]

For $r/\rho$ sufficiently small, $v\ne0$.
Moreover, the elementary normalization inequality
\[
\left\|
\frac{v}{\|v\|_2}-u
\right\|_2
\le
\frac{2\|v-au\|_2}{a}
\]
yields
\[
\left\|
\widehat n(p,r)-|n(p)|
\right\|_2
\le
\frac{2\sqrt d\,C_d^{\mathrm{stab}}}
{\omega_{d-1}}
\frac r\rho.
\]
This proves the claim.
\end{proof}


\begin{corollary}[Normal estimation from a cubical discretization]
\label{cor:cubical-normal-estimation}
Under the assumptions of
Corollary~\ref{cor:local-normal-estimation},
let $U_h(X)$ be the union of all closed
axis-aligned grid cubes of side length $h$
that intersect $X$, and define
\[
Y_h(p,r)=U_h(X)\cap B(p,r).
\]

For each $i\in\{1,\ldots,d\}$, set
\[
V_{i,h}(p,r)
:=
\mathcal H^{d-1}
\bigl(\pi_{\widehat i}(Y_h(p,r))\bigr),
\]
and
\[
\widehat n_{i,h}(p,r)
:=
\frac{V_{i,h}(p,r)}
{\left(\sum_{j=1}^d V_{j,h}(p,r)^2\right)^{1/2}}.
\]

There exist constants $c_d,C_d>0$, depending
only on $d$, such that, whenever
\[
0<h\le r,\qquad
\frac hr+\frac r\rho<c_d,
\]
the estimator is well-defined and satisfies
\[
\boxed{
\left\|
\widehat n_h(p,r)-|n(p)|
\right\|_2
\le
C_d\left(\frac hr+\frac r\rho\right).
}
\]
\end{corollary}

\begin{proof}
By the local stability theorem for cubical
discretizations, applied with $k=m=d-1$,
there exists a constant
$C_d^{\mathrm{stab}}>0$ such that
\[
\left|
V_{i,h}(p,r)
-
\mathcal H^{d-1}
\bigl(\pi_{\widehat i}(B_T(p,r))\bigr)
\right|
\le
C_d^{\mathrm{stab}}r^{d-2}
\left(h+\frac{r^2}{\rho}\right),
\]
where $T=T_pX$.

Using again the exact projection formula
\[
\mathcal H^{d-1}
\bigl(\pi_{\widehat i}(B_T(p,r))\bigr)
=
\omega_{d-1}r^{d-1}|n_i(p)|,
\]
we obtain
\[
\left|
V_{i,h}(p,r)
-
\omega_{d-1}r^{d-1}|n_i(p)|
\right|
\le
C_d^{\mathrm{stab}}r^{d-2}
\left(h+\frac{r^2}{\rho}\right).
\]

Writing
\[
v_h=(V_{1,h}(p,r),\ldots,V_{d,h}(p,r)),
\qquad
a=\omega_{d-1}r^{d-1},
\qquad
u=|n(p)|,
\]
gives
\[
\|v_h-au\|_2
\le
\sqrt d\,C_d^{\mathrm{stab}}r^{d-2}
\left(h+\frac{r^2}{\rho}\right).
\]

For $h/r+r/\rho$ sufficiently small,
$v_h\ne0$. Applying the normalization
inequality yields
\[
\left\|
\widehat n_h(p,r)-|n(p)|
\right\|_2
\le
\frac{2\sqrt d\,C_d^{\mathrm{stab}}}
{\omega_{d-1}}
\left(\frac hr+\frac r\rho\right),
\]
which proves the result.
\end{proof}


\begin{remark}[Sign ambiguity]
The projection volumes determine only the
absolute values of the normal components.
In particular, the relative signs of
$n_1(p),\ldots,n_d(p)$ cannot be recovered
from these measurements alone.
Additional geometric information is needed
to reconstruct the full normal direction.
\end{remark}
